\section{Concluding Remarks}
\label{sec_conc_remarks}

The composite hypothesis of Section~\ref{sec_intro_hypothesis} is supported across all three clauses,
subject to the qualifications set out in Section~\ref{sec_conc_verdict}.
The experiments show that the quality of the per-keypoint embedding,
rather than the complexity of the grouping algorithm, is the main constraint on grouping performance.
PG-GAT demonstrates that a skeleton-aware graph attention network can learn such an embedding through
synthetic pre-training followed by COCO fine-tuning,
although its performance remains below the jointly trained reference and the claim of competitive performance
is therefore qualitative.
The same representation also supports autonomous person-count estimation and reuse across multiple detectors,
with the detector-swap experiments providing stronger evidence for the third clause than originally required
by the hypothesis.
The resulting contribution is a grouping module that can be developed, evaluated,
and replaced independently of the detector,
at a measured cost of approximately $0.02$ PGA relative to native grouping across the three detector families tested.
Across the wider investigation, one result remained consistent:
improvements to the per-keypoint embedding produced larger gains in grouping quality than changes to the
grouping algorithm applied to it.
